\subsubsection{Planteamiento del problema}

La incorporación de herramientas digitales en la gestión de la información en salud ha transformado profundamente la forma en que las instituciones sanitarias organizan y comparten los datos de sus pacientes. En las últimas dos décadas, el uso de registros electrónicos ha crecido de manera considerable en distintos países, impulsado por la necesidad de contar con información oportuna, completa y accesible \cite{gurupur2024standards}. Sin embargo, este avance no ha sido uniforme: mientras algunos sistemas hospitalarios de alta complejidad han avanzado en la digitalización de sus procesos, los centros de atención de menor escala, y en particular las clínicas odontológicas privadas, siguen operando con flujos de información fragmentados, dispersos entre herramientas que no se comunican entre sí. Este desfase no es un problema menor, pues incide directamente en la velocidad y la calidad con que el personal puede atender a cada paciente.

A nivel mundial, la literatura coincide en señalar que la información clínica dispersa y sin estandarización genera una carga de trabajo adicional sobre el personal de salud. Cuando los datos de un paciente están repartidos entre sistemas distintos, sin conexión entre ellos, el personal dedica tiempo considerable a trasladar información de forma manual de un registro a otro, a corregir errores de captura y a buscar datos que deberían estar disponibles de inmediato \cite{gurupur2024standards}. Esta situación se traduce en que una porción importante del tiempo laboral en los centros de salud se destina a tareas administrativas y documentales en lugar de a la atención directa del paciente \cite{murad2024documentation}. En regiones como América Latina, este problema adquiere dimensiones adicionales: un análisis comparativo del marco normativo de 26 países de la región, incluido Ecuador, identificó que la falta de sistemas de información bien diseñados e interoperables mantiene silos de información que generan ineficiencias y dificultan el acceso e intercambio de datos entre los actores involucrados en la atención \cite{bagolle2020bid}. Esta problemática, lejos de ser exclusiva de los grandes hospitales, afecta de igual manera a los centros de salud de menor tamaño que no cuentan con recursos ni estructura para integrar sus procesos.

En el sector odontológico en particular, la evidencia muestra que la dependencia de registros en papel y procesos manuales provoca retrasos medibles y riesgos de omisión en la información clínica. Un estudio comparativo realizado en una clínica dental encontró que identificar antecedentes relevantes en registros físicos de pacientes tomaba aproximadamente dos minutos por expediente, mientras que el tiempo total para registrar y generar reportes de cien pacientes ascendió a más de trescientos minutos usando papel, frente a cuarenta minutos con un software especializado; además, con registros en papel se pasaron por alto condiciones médicas críticas que podían haber comprometido la seguridad del paciente \cite{scarano2025digital}. En el mismo sentido, la experiencia de implementación de un sistema de registro electrónico en clínicas dentales públicas documentó que los usuarios reportaron procesos excesivamente lentos y una carga de trabajo adicional derivada de la falta de adaptación del sistema a las necesidades específicas del entorno odontológico \cite{hajisidek2017ehr}. En Ecuador, una encuesta realizada a más de trescientos odontólogos reveló que el principal temor frente a la adopción de nuevas herramientas digitales es la inseguridad respecto al manejo de los datos de los pacientes, seguido por la falta de tiempo para aprender a usar sistemas nuevos \cite{cherrezojeda2020ict}; esto evidencia que las barreras no son solo tecnológicas, sino también organizacionales y culturales. En las clínicas odontológicas privadas de tamaño pequeño o mediano, donde no existe un área de tecnología ni procesos estandarizados, estas dificultades se intensifican y terminan recayendo sobre el personal que atiende a los pacientes.

El Centro Odontológico Bellidental, ubicado en el cantón Baños, provincia de Tungurahua, refleja de manera puntual la situación descrita. Con una atención de aproximadamente setenta y cinco pacientes por semana y más de mil expedientes acumulados, el centro opera con un modelo completamente manual y desconectado: las historias clínicas se registran en hojas de cálculo que replican el formato oficial del Ministerio de Salud Pública, las citas se anotan en un calendario compartido en línea sin integración con los datos del paciente, los pagos se documentan en hojas de papel firmadas a mano, las recetas se escriben a mano y los recordatorios a los pacientes se envían uno a uno por mensajes de texto a través de una aplicación de mensajería instantánea. Cada una de estas herramientas funciona de manera independiente, sin que exista un punto central donde la información converja. En la práctica, esto obliga al personal a interrumpir constantemente las actividades clínicas para buscar datos entre filas de cálculo, transcribir información entre distintos registros y coordinar mediante mensajes individuales cada gestión administrativa. El tiempo que se pierde en estas tareas de oficina es directamente tiempo que no se dedica a la atención del paciente.

La situación del Centro Odontológico Bellidental tiene raíces estructurales claras: la ausencia de un flujo de información clínico-administrativa estandarizado, el uso de registros dispersos sin integración entre sí y la dependencia de canales de comunicación no articulados entre las distintas etapas del proceso de atención. Si esta condición se mantiene sin cambios, es previsible que el tiempo dedicado a tareas de gestión por cada paciente continúe creciendo a medida que aumenta el volumen de atenciones, que la saturación del personal administrativo y clínico se agudice, que los errores por transcripción manual y las omisiones de información se acumulen, y que la capacidad del centro para mantener la calidad del servicio y competir con otros establecimientos de salud se vea progresivamente comprometida. Ante este escenario, surge la siguiente pregunta de investigación: ¿De qué manera el nivel de estandarización e integración del flujo de información clínico-administrativa incide en el tiempo promedio de gestión administrativa por paciente en el Centro Odontológico Bellidental?